\section{Introduction}

Looped (weight-tied) Transformers reuse a single block $f$ for $N$ iterations ($h \leftarrow h + f(h)$, same $f$ at each step), increasing effective depth without adding parameters.
This design appears in Universal Transformers~\citep{dehghani2018universal}, ALBERT~\citep{lan2019albert}, and recent work on algorithmic reasoning and latent computation~\citep{yang2023looped,fan2024looped,saunshi2025latent,gatmiry2024multistep,ng2024loopnn,zhu2025ouro}.

\begin{figure*}[t]
  \centering
  \resizebox{\textwidth}{!}{\input{figures/main_schematic.tikz}}
  \caption{\textbf{Weight sharing changes how residuals accumulate.}
  \textbf{(a)}~In a deep network with independent weights, residual updates point in different directions and accumulate like a random walk, with norm $\Theta(\!\sqrt{L})$. Standard scaling $\varepsilon=1/\!\sqrt{L}$ keeps the output bounded.
  \textbf{(b)}~In a looped network, a single block is reused $N$ times. The shared weights make successive updates align, so their sum grows as $\Theta(N)$---requiring the stronger scaling $\varepsilon=1/N$.}
  \label{fig:main-schematic}
\vspace{-0.3cm}
\end{figure*}


In practice, increasing $N$ often leads to training instability such as exploding hidden states and high sensitivity to the learning rate~\citep{zhu2025ouro}.
A standard remedy is to scale each residual branch by a factor $\varepsilon$ that shrinks with depth, giving $h \leftarrow h + \varepsilon\,f(h)$.
Prior depth-scaling analyses prescribe $\varepsilon = 1/\!\sqrt{N}$ for deep residual networks~\citep{bordelon2024depthwise,dey2025completep}, but whether this rule transfers to looped architectures---where the same $f$ is reused at every step---has not been established.

We find that $1/\!\sqrt{N}$ scaling is indeed insufficient for looped models.
Consider the residual-stream norm $\|h_N\|$ after $N$ iterations.
For standard (non-shared) deep networks, $\varepsilon = 1/\!\sqrt{L}$ successfully keeps $\|h_L\|$ bounded as depth $L$ grows (Figure~\ref{fig:energy}, top row).
For looped networks, however, $\varepsilon = 1/\!\sqrt{N}$ fails to control $\|h_N\|$, which grows rapidly with $N$; in contrast, $\varepsilon = 1/N$ keeps it bounded (Figure~\ref{fig:energy}, bottom row).
Our theoretical analysis (Section~\ref{sec:theory}) explains this discrepancy: the $1/\!\sqrt{N}$ rule relies on the assumption that each layer has independent weights, but weight sharing makes successive updates correlated, amplifying residual-stream norm growth from $\Theta(\!\sqrt{N})$ to $\Theta(N)$.

\begin{figure*}[t]
  \centering
  \includegraphics[width=\textwidth]{figures/variance.pdf}
  \caption{\textbf{Linear scaling stabilizes looped networks; sqrt scaling does not.}
  Normalized residual-stream norm $R=d^{-1/2}\|h\|_2$ (log scale) vs.\ depth $L$ (top) or loop count $N$ (bottom) during the first 10 training steps in the Llama-style pre-norm Transformer diagnostic; lines colored by step.
  $1/\!\sqrt{L}$ scaling stabilizes deep networks (panel~b), but $1/\!\sqrt{N}$ fails for loops (panel~e).
  Linear scaling $\varepsilon{=}1/N$ keeps the residual-stream norm bounded across loop counts (panel~f).}
  \label{fig:energy}
\vspace{-0.3cm}
\end{figure*}


Beyond stabilizing the forward pass, the $1/N$ scaling also fixes the learning rate.
The same constructive accumulation that drives $\Theta(N^2)$ squared-norm growth also amplifies weight updates: the output change from one optimizer step scales as $\eta\varepsilon N$, so setting $\varepsilon=1/N$ makes the stable learning rate constant in $N$ (Section~\ref{sec:theory}).
This enables hyperparameter transfer: a learning rate tuned at $N{=}1$ remains near-optimal at larger $N$ without re-tuning.

We further extend the analysis to practical multi-layer blocks, where $L$ distinct layers are each reused $N$ times (Section~\ref{sec:multilayer}).
This introduces a second variance source: across independent layers, updates accumulate as a random walk in $L$, just as in standard deep networks.
The factorized parameterization $\varepsilon = \lambda/(N\sqrt{L})$ handles both sources independently: $1/N$ cancels the within-layer quadratic growth, while $1/\!\sqrt{L}$ controls the across-layer random walk.
The resulting learning-rate law, $\eta \lesssim 1/(\lambda\sqrt{L})$, depends only on the unique depth $L$, not on $N$, so hyperparameter transfer continues to hold for multi-layer blocks.

Experiments on decoder-only Transformers trained on FineWeb-Edu~\citep{penedo2024fineweb} confirm these predictions (Section~\ref{sec:experiments}).
The pairwise correlation structure underlying quadratic accumulation---dense positive cosine similarity between loop-step updates---persists well beyond initialization through full training, confirming that the $\Theta(N^2)$ growth mechanism is not an initialization artifact.
Linear residual scaling yields improved trainability and consistent learning-rate transfer across $N\in\{1,2,4,8\}$: the optimal learning rate remains nearly invariant, unlike under $1/\!\sqrt{N}$ scaling where it shifts with $N$.
The factorized parameterization $\varepsilon=\lambda/(N\sqrt{L})$ further extends this transfer across depths $L\in\{12,24,48\}$, with a single learning rate remaining near-optimal over all tested $(N,L)$ combinations.

Taken together, our work extends the depth-scaling framework~\citep{yang2022tensorprogramsV,dey2025completep} to weight-shared architectures, showing that the standard independence assumptions break down under parameter reuse and deriving the corrected scaling rules.
Beyond the theoretical contribution, the resulting parameterization directly solves two practical problems: it improves trainability at large loop counts and eliminates the need to re-tune hyperparameters when varying $N$.
By making the loop count stable and tunable, these results establish reuse as a practical scaling axis for weight-tied Transformers.
